% GENERATED FILE. Edit canonical lessons, then run npm run generate:books.
% canonical-source-hash: fnv1a64:e3e08d2c3aea2464
% canonical-lessons: PT-C279-imprensa, PT-C279-reportagem, PT-C279-reporter, PT-C279-correspondente, PT-C279-manchete

\chapter{A imprensa}
\label{ch:pt-a-imprensa}
\hlchaptermodality{279}

\begin{chapteropening}
\textbf{By the end of this chapter:} \emph{I can talk about the press: reports, reporters, correspondents and headlines.}

Complete the last lesson of chapter 279: I can talk about the press: reports, reporters, correspondents and headlines.
\end{chapteropening}

\section[a imprensa]{a imprensa — the press, the newspapers}

\label{lesson:PT-C279-imprensa}



\begin{quote}
Before the new one: say the Portuguese for live, then the Portuguese for the broadcast.
\end{quote}

\subsection*{You'll want to know: a imprensa}
\textbf{a imprensa} — \textquotedblleft{}the press, the newspapers\textquotedblright{}.

\begin{quote}
A imprensa falou muito do assunto.
\end{quote}

\begin{quote}
\emph{The press talked a lot about the matter.}
\end{quote}

It also names the printing press.

\begin{cousinweb}
\emph{imprensa} goes back to Latin \emph{impressus} \textquotedblleft{}pressed in, printed\textquotedblright{} (EtymDB), from \emph{imprimere} \textquotedblleft{}to press into\textquotedblright{}, built on \emph{premere} \textquotedblleft{}to press\textquotedblright{}.
\end{cousinweb}

\begin{grammarlens}[title={What you've built}]
One more word for telling, explaining and describing at B1.
\end{grammarlens}

\subsection*{Your turn}
\begin{itemize}
\raggedright
  \item \emph{Say it:} \emph{a imprensa}
  \item \emph{Say it:} \emph{a imprensa}, once more
  \item \emph{Say it:} a sentence of your own with \emph{a imprensa}
  \item \emph{Recall:} say the Portuguese for live, then the Portuguese for the broadcast, then say \emph{a imprensa} again
\end{itemize}

\begin{tcolorbox}[breakable,colback=teal!4,colframe=teal!35!black,title={Before you move on}]
What does a imprensa mean? (\textquotedblleft{}The press, the newspapers\textquotedblright{}.) Say it once more.
\end{tcolorbox}
\section[a reportagem]{a reportagem — the report, the feature story}

\label{lesson:PT-C279-reportagem}



\begin{quote}
Before the new one: say the Portuguese for the broadcast, then the Portuguese for the press.
\end{quote}

\subsection*{You'll want to know: a reportagem}
\textbf{a reportagem} — \textquotedblleft{}the report, the feature story\textquotedblright{}.

\begin{quote}
Vi uma reportagem sobre a vida no deserto.
\end{quote}

\begin{quote}
\emph{I saw a report about life in the desert.}
\end{quote}

A longer piece on TV or in a paper, with interviews and pictures.

\begin{cousinweb}
\emph{reportagem} came from French \emph{reportage} (EtymDB), built on \emph{reporter} \textquotedblleft{}to report\textquotedblright{}, which goes back to Latin \emph{reportāre} \textquotedblleft{}to carry back\textquotedblright{}, from \emph{portāre} \textquotedblleft{}to carry\textquotedblright{}.
\end{cousinweb}

\begin{grammarlens}[title={What you've built}]
One more word for telling, explaining and describing at B1.
\end{grammarlens}

\subsection*{Your turn}
\begin{itemize}
\raggedright
  \item \emph{Say it:} \emph{a reportagem}
  \item \emph{Say it:} \emph{a reportagem}, once more
  \item \emph{Say it:} a sentence of your own with \emph{a reportagem}
  \item \emph{Recall:} say the Portuguese for the broadcast, then the Portuguese for the press, then say \emph{a reportagem} again
\end{itemize}

\begin{tcolorbox}[breakable,colback=teal!4,colframe=teal!35!black,title={Before you move on}]
What does a reportagem mean? (\textquotedblleft{}The report, the feature story\textquotedblright{}.) Say it once more.
\end{tcolorbox}
\section[o repórter]{o repórter — the reporter}

\label{lesson:PT-C279-reporter}



\begin{quote}
Before the new one: say the Portuguese for the press, then the Portuguese for the report.
\end{quote}

\subsection*{You'll want to know: o repórter}
\textbf{o repórter} — \textquotedblleft{}the reporter\textquotedblright{}.

\begin{quote}
Um repórter falou com os moradores.
\end{quote}

\begin{quote}
\emph{A reporter spoke to the residents.}
\end{quote}

The same form for a woman: \emph{a repórter}.

\begin{cousinweb}
\emph{repórter} was borrowed from English \emph{reporter} (EtymDB), a word that goes back to Latin \emph{reportāre} \textquotedblleft{}to carry back\textquotedblright{}, from \emph{portāre} \textquotedblleft{}to carry\textquotedblright{}.
\end{cousinweb}

\begin{grammarlens}[title={What you've built}]
One more word for telling, explaining and describing at B1.
\end{grammarlens}

\subsection*{Your turn}
\begin{itemize}
\raggedright
  \item \emph{Say it:} \emph{o repórter}
  \item \emph{Say it:} \emph{o repórter}, once more
  \item \emph{Say it:} a sentence of your own with \emph{o repórter}
  \item \emph{Recall:} say the Portuguese for the press, then the Portuguese for the report, then say \emph{o repórter} again
\end{itemize}

\begin{tcolorbox}[breakable,colback=teal!4,colframe=teal!35!black,title={Before you move on}]
What does o repórter mean? (\textquotedblleft{}The reporter\textquotedblright{}.) Say it once more.
\end{tcolorbox}
\section[o correspondente]{o correspondente — the correspondent}

\label{lesson:PT-C279-correspondente}



\begin{quote}
Before the new one: say the Portuguese for the report, then the Portuguese for the reporter.
\end{quote}

\subsection*{You'll want to know: o correspondente}
\textbf{o correspondente} — \textquotedblleft{}the correspondent\textquotedblright{}.

\begin{quote}
O correspondente do jornal vive no estrangeiro.
\end{quote}

\begin{quote}
\emph{The newspaper’s correspondent lives abroad.}
\end{quote}

As an adjective it means ‘matching’: \emph{a resposta correspondente}.

\begin{cousinweb}
\emph{correspondente} is a learned word from Medieval Latin \emph{correspondēns} \textquotedblleft{}answering to each other\textquotedblright{}, built on \emph{respondēre} \textquotedblleft{}to answer\textquotedblright{} (Lewis \& Short).
\end{cousinweb}

\begin{grammarlens}[title={What you've built}]
One more word for telling, explaining and describing at B1.
\end{grammarlens}

\subsection*{Your turn}
\begin{itemize}
\raggedright
  \item \emph{Say it:} \emph{o correspondente}
  \item \emph{Say it:} \emph{o correspondente}, once more
  \item \emph{Say it:} a sentence of your own with \emph{o correspondente}
  \item \emph{Recall:} say the Portuguese for the report, then the Portuguese for the reporter, then say \emph{o correspondente} again
\end{itemize}

\begin{tcolorbox}[breakable,colback=teal!4,colframe=teal!35!black,title={Before you move on}]
What does o correspondente mean? (\textquotedblleft{}The correspondent\textquotedblright{}.) Say it once more.
\end{tcolorbox}
\section[a manchete]{a manchete — the headline}

\label{lesson:PT-C279-manchete}



\begin{quote}
Before the new one: say the Portuguese for the reporter, then the Portuguese for the correspondent.
\end{quote}

\subsection*{You'll want to know: a manchete}
\textbf{a manchete} — \textquotedblleft{}the headline\textquotedblright{}.

\begin{quote}
A manchete de hoje fala do governo.
\end{quote}

\begin{quote}
\emph{Today’s headline is about the government.}
\end{quote}

The big headline on the front page; \emph{o título} is any heading.

\begin{cousinweb}
\emph{manchete} was borrowed from French \emph{manchette} \textquotedblleft{}a cuff, a headline\textquotedblright{}, a little \emph{manche} \textquotedblleft{}sleeve\textquotedblright{}, which comes from Latin \emph{manica} \textquotedblleft{}a long sleeve\textquotedblright{}, built on \emph{manus} \textquotedblleft{}hand\textquotedblright{} (Lewis \& Short).
\end{cousinweb}

\begin{grammarlens}[title={What you've built}]
One more word for telling, explaining and describing at B1.
\end{grammarlens}

\subsection*{Your turn}
\begin{itemize}
\raggedright
  \item \emph{Say it:} \emph{a manchete}
  \item \emph{Say it:} \emph{a manchete}, once more
  \item \emph{Say it:} a sentence of your own with \emph{a manchete}
  \item \emph{Recall:} say the Portuguese for the reporter, then the Portuguese for the correspondent, then say \emph{a manchete} again
\end{itemize}

\begin{tcolorbox}[breakable,colback=teal!4,colframe=teal!35!black,title={Before you move on}]
What does a manchete mean? (\textquotedblleft{}The headline\textquotedblright{}.) Say it once more.
\end{tcolorbox}
